%% Abstract should be no more than 250 words
% This study presents a rigorous verification effort for a computational model of a simplified medical device geometry involving conjugate heat transfer. Using the open-source \texttt{chtMultiRegionFoam} solver in OpenFOAM, both code and solution verification activities were performed to establish the model’s numerical accuracy and trustworthiness prior to validation. Code verification was carried out using benchmarks: Hagen–Poiseuille flow for laminar momentum transport and solid conduction. In both cases, the numerical results closely matched analytical solutions. Solution verification included systematic studies of mesh resolution, timestep size, and solver tolerance. Spatial and temporal errors were quantified using an extrapolated value computed using StREEQ to generate relative error measures, with observed convergence behavior supporting the selection of a configuration that balances computational cost and accuracy. The modeling effort aligns with the ASME V\&V 40 standard and follows a risk-informed credibility framework suitable for medium to low risk applications. This work provides a transparent and reproducible case study in model verification for computational modeling and simulation 
% (CM\&S) in medical device applications. The results provide a foundation for subsequent validation and uncertainty quantification, contributing to the broader goal of establishing credible simulation-based evidence.

Computational models and simulations (CM\&S) are increasingly used in medical device development, yet their adoption in regulatory decision-making remains limited due to challenges in establishing model credibility. This study presents a structured and reproducible verification framework for a computational model of a simplified medical device involving conjugate heat transfer using \texttt{OpenFOAM}'s \texttt{chtMultiRegionFoam} solver. Code verification is performed using analytical benchmarks for laminar pipe flow, radial heat conduction, and transient thermal diffusion yielding observed orders of convergence of $p=1.95$, $p=2.08$, and $p=2.15$ (at 0.025 s) respectively. Solution verification evaluates discretization error through systematic mesh and timestep refinement and regression-based extrapolation using the \texttt{StREEQ} method. Results show that spatial and temporal discretization errors in transient coupled simulations interact and cannot be assessed independently. Co-refinement is required to capture true convergence behavior, while independent refinement yields polluted convergence rates. Computational cost analysis demonstrates diminishing returns with refinement and highlights that runtime does not vary monotonically with discretization parameters due to solver convergence behavior. The use of a simplified geometry enables rigorous verification while retaining the dominant physics of the system. Grounded in a risk-informed framework consistent with ASME V\&V 40 Standard, this study demonstrates that a level of verification rigor beyond minimal practice can be achieved through targeted model simplification and systematic analysis under realistic computational constraints.